% Generated draft: requires mathematical and source review.
\documentclass[11pt]{article}
\usepackage[margin=25mm]{geometry}
\usepackage{fontspec}
\setmainfont{Latin Modern Roman}
\usepackage{amsmath,amssymb,mathtools,unicode-math}
\setmathfont{Latin Modern Math}
\usepackage{enumitem,needspace,xurl,hyperref}
\setlength{\parindent}{0pt}
\setlength{\parskip}{5pt}
\setlength{\emergencystretch}{3em}
\newcommand{\N}{\mathbb{N}}
\newcommand{\Z}{\mathbb{Z}}
\newcommand{\Q}{\mathbb{Q}}
\newcommand{\R}{\mathbb{R}}
\newcommand{\C}{\mathbb{C}}
\newcommand{\F}{\mathbb{F}}
\newcommand{\A}{\mathbb{A}}
\newcommand{\PP}{\mathbb{P}}
\newcommand{\E}{\mathbb{E}}
\newcommand{\eps}{\varepsilon}
\newcommand{\dd}{\,\mathrm{d}}
\begin{document}
\section*{UnsolvedMath ID: 30005426}
\textbf{Source problem code:} catalogue entry
\subsection{Definitions and mathematical statement}
The supplied question concerns representation-finite algebras in the context of finite-dimensional algebras. Its ground field, grading conventions, and precise meaning of Koszul dual are not specified. The following definitions describe a conditional formulation; these missing qualifications require review.

Let \(k\) be a field and let \(B\) be a finite-dimensional \(k\)-algebra. Write \(\operatorname{mod}B\) for the category of finite-dimensional right \(B\)-modules. The algebra \(B\) is \emph{representation-finite} if this category has only finitely many isomorphism classes of indecomposable objects. Choose representatives \(M_1,\ldots,M_t\) of these classes and put
\[
M=\bigoplus_{i=1}^{t}M_i,\qquad A=\operatorname{End}_B(M).
\]
With this module-side convention, \(A\) is the Auslander algebra of \(B\).

For a finite-dimensional algebra \(C\), its right global dimension \(\operatorname{gldim}C\) is the supremum of the projective dimensions of right \(C\)-modules. Its right dominant dimension \(\operatorname{domdim}C\) is the supremum of the integers \(r\geq 0\) such that \(I^0,\ldots,I^{r-1}\) are projective in a minimal injective resolution
\[
0\longrightarrow C_C\longrightarrow I^0\longrightarrow I^1\longrightarrow\cdots.
\]
The supremum may be infinite. Under the convention used here, \(C\) is \emph{higher Auslander} if there is an integer \(n\geq 1\) such that
\[
\operatorname{gldim}C\leq n+1\leq\operatorname{domdim}C.
\]
The intended indexing convention and treatment of semisimple algebras should be confirmed from the source.

To specify one usual meaning of Koszul dual, suppose additionally that \(A\) has a nonnegative grading \(A=\bigoplus_{j\geq0}A_j\), with \(A_0\) semisimple, and is Koszul: the graded right module \(A_0=A/A_{>0}\) has a graded projective resolution whose term in homological degree \(i\) is generated in internal degree \(i\). One convention then defines
\[
A^{!}=\bigoplus_{i\geq0}\operatorname{Ext}^{i}_A(A_0,A_0)
\]
with Yoneda multiplication, using Ext over the underlying algebra and grading the result by homological degree. Opposite-algebra conventions also occur. The supplied record neither gives such a grading nor asserts that every Auslander algebra under consideration is Koszul. It also does not say whether an Ext algebra without a Koszul hypothesis is intended.

\textbf{Classification question.} Determine which representation-finite algebras \(B\) have the property that the intended Koszul dual \(A^{!}\) of their Auslander algebra is a finite-dimensional higher Auslander algebra; with the convention above, this means
\[
\exists n\in\mathbb Z_{\geq1}\quad
\operatorname{gldim}(A^{!})\leq n+1\leq\operatorname{domdim}(A^{!}).
\]
This formulation is conditional on resolving the field, grading, duality, and higher Auslander conventions. It does not impose a Koszul hypothesis on the original question without source confirmation. The supplied registry labels the question open; this is a supplied status, not an independently verified mathematical assessment.
\subsection{Short English statement}
Classify the representation-finite algebras whose Auslander algebras have higher Auslander Koszul duals. The supplied record leaves the ground field and the meaning and hypotheses of Koszul duality unspecified, so the displayed formulation remains conditional and requires human mathematical review.
\subsection{Sources}
\begin{itemize}
\item Local UnsolvedMath registry: supplied complete record with permanent ID \(30005426\), source problem code \texttt{catalogue entry}, title \emph{Higher-Auslander Koszul Duals of Auslander Algebras}.
\item \emph{Representation Theory of Quivers and Finite-Dimensional Algebras} (2023), Oberwolfach Reports, pp.~397--486, DOI \texttt{10.4171/owr/2023/7}. Supplied source report and literature-source link: \url{https://doi.org/10.4171/owr/2023/7}.
\item The supplied original statement attributes the contribution \emph{Auslander regular algebras and Koszul duality} to Ren\'e Marczinzik, joint work with Aaron Chan and Osamu Iyama. No separate proposer is specified; the registry gives proposed year 2023.
\item Chan--Iyama--Marczinzik, \emph{Dominant Auslander--Gorenstein algebras}: \url{https://arxiv.org/abs/2210.06180}. This is the additional supplied literature-source link.
\item Literature assessment dated exactly \texttt{2026-08-22}: the supplied curation labels the question open and reports that the cited work establishes important hereditary and self-injective cases without a general classification. This is supplied dated curation, not independently verified here. The background explicitly requests independent verification of time-sensitive status and cited claims.
\item The registry labels the statement \texttt{corrected\_verified} and reports comparison with the official Oberwolfach report. That verification claim is supplied provenance; no source visit or independent verification is asserted in this draft.
\end{itemize}
\end{document}
